\frfilename{mt225.tex}
\versiondate{16.8.15}
\copyrightdate{1996}

\def\chaptername{Teorema Dasar Kalkulus}
\def\sectionname{Fungsi kontinu mutlak}

\newsection{225}

Sekarang kita siap memberikan pencirian lengkap fungsi-fungsi yang dapat
muncul sebagai integral tak tentu (225E, 225Xf). Gagasan pokoknya ialah
`kontinuitas mutlak' (225B). Pada paruh kedua bagian ini (225G-225N), saya
menguraikan beberapa hubungan antara konsep tersebut dan konsep-konsep yang
telah kita jumpai.

\leader{225A}{Kontinuitas mutlak integral tak
\dvrocolon{tentu}}\cmmnt{ Saya memulai dengan suatu hasil dasar yang mudah dari teori
ukuran umum.

\medskip

\noindent}{\bf Teorema} Misalkan $(X,\Sigma,\mu)$ sembarang ruang ukur dan
$f$ suatu fungsi bernilai real terintegralkan yang didefinisikan pada suatu
subhimpunan koterabaikan dari $X$. Maka untuk setiap $\epsilon>0$ terdapat
suatu himpunan terukur $E$ berukuran berhingga dan suatu bilangan real
$\delta>0$ sedemikian sehingga $\int_F|f|\le\epsilon$ kapan pun
$F\in\Sigma$ dan $\mu(F\cap E)\le\delta$.

\proof{ Terdapat suatu barisan tak-menurun $\sequencen{g_n}$ dari fungsi
sederhana nonnegatif sedemikian sehingga
$|f|\eae\lim_{n\to\infty}g_n$ dan
$\int|f|=\lim_{n\to\infty}\int g_n$. Ambil $n\in\Bbb N$ sedemikian sehingga
$\int g_n\ge\int|f|-\bover12\epsilon$.
Ambil $M>0$ dan $E\in\Sigma$ sedemikian sehingga $\mu E<\infty$ dan
$g_n\le M\chi E$; tetapkan $\delta=\epsilon/2M$.
Jika $F\in\Sigma$ dan $\mu(F\cap E)\le\delta$, maka

\Centerline{$\int_Fg_n=\int g_n\times\chi F
\le M\mu(F\cap E)\le\Bover12\epsilon$;}

\noindent akibatnya

\Centerline{$\int_F|f|
=\int_Fg_n+\int_F|f|-g_n\le\Bover12\epsilon+\int|f|-g_n\le\epsilon$.}
}%end of proof of 225A

\leader{225B}{Fungsi kontinu mutlak pada $\Bbb R$: Definisi}
Jika $[a,b]$ suatu interval tertutup tak kosong dalam $\Bbb R$ dan
$f:[a,b]\to\Bbb R$ suatu fungsi, kita mengatakan bahwa $f$ {\bf kontinu
mutlak} jika untuk setiap $\epsilon>0$ terdapat $\delta>0$ sedemikian sehingga
$\sum_{i=1}^n|f(b_i)-f(a_i)|\le\epsilon$ kapan pun
$a\le a_1\le b_1\le a_2\le b_2\le\ldots\le a_n\le b_n\le b$ dan
$\sum_{i=1}^nb_i-a_i\le\delta$.

\cmmnt{\medskip

\noindent{\bf Catatan} Ungkapan `kontinu mutlak' digunakan dalam berbagai
pengertian dalam teori ukuran. Pengertian-pengertian itu berkaitan erat (jika
dilihat dengan cara yang tepat), tetapi tidak identik; Anda perlu senantiasa
mencermati konteks setiap definisi.
}%end of comment

\leader{225C}{Proposisi} Misalkan $[a,b]$ suatu interval tertutup tak kosong
dalam $\Bbb R$.

(a) Jika $f:[a,b]\to\Bbb R$ kontinu mutlak, maka fungsi tersebut kontinu seragam.

(b) Jika $f:[a,b]\to\Bbb R$ kontinu mutlak, maka fungsi tersebut bervariasi terbatas
pada $[a,b]$, sehingga terdiferensialkan hampir di mana-mana dalam
$[a,b]$, dan turunannya terintegralkan pada $[a,b]$.

(c) Jika $f$, $g:[a,b]\to\Bbb R$ kontinu mutlak, maka demikian pula $f+g$
dan $cf$, untuk setiap $c\in\Bbb R$.

(d) Jika $f$, $g:[a,b]\to\Bbb R$ kontinu mutlak, maka demikian pula
$f\times g$.

(e) Jika $g:[a,b]\to [c,d]$ dan $f:[c,d]\to\Bbb R$ kontinu mutlak, serta
$g$ tak-menurun, maka komposisi $fg:[a,b]\to\Bbb R$ kontinu mutlak.

\proof{{\bf (a)} Misalkan $\epsilon>0$. Maka terdapat $\delta>0$ sedemikian
sehingga $\sum_{i=1}^n|f(b_i)-f(a_i)|\le\epsilon$ kapan pun
$a\le a_1\le b_1\le a_2\le b_2\le\ldots\le a_n\le b_n\le b$ dan
$\sum_{i=1}^nb_i-a_i\le\delta$; tetapi tentu saja sekarang
$|f(y)-f(x)|\le\epsilon$ kapan pun $x$, $y\in[a,b]$ dan
$|x-y|\le\delta$. Karena $\epsilon$ sebarang, $f$ kontinu seragam.



\medskip

{\bf (b)} Misalkan $\delta>0$ sedemikian sehingga
$\sum_{i=1}^n|f(b_i)-f(a_i)|\le 1$ kapan pun
$a\le a_1\le b_1\le a_2\le b_2\le\ldots\le a_n\le b_n\le b$ dan
$\sum_{i=1}^nb_i-a_i\le\delta$. Jika
$a\le c=c_0\le c_1\le\ldots\le c_n\le d\le\min(b,c+\delta)$, maka
$\sum_{i=1}^n|f(c_i)-f(c_{i-1})|\le 1$, sehingga
$\Var_{[c,d]}(f)\le 1$; oleh karena itu (dengan induksi terhadap $k$ dan
menggunakan 224Cc untuk langkah induksi),
$\Var_{[a,\min(a+k\delta,b)]}(f)\le k$ untuk setiap $k$, dan

\Centerline{$\Var_{[a,b]}(f)\le\lceil(b-a)/\delta\rceil<\infty$.}

Karena itu, menurut 224I, $f'$ terintegralkan.

\medskip

{\bf (c)(i)} Misalkan $\epsilon>0$. Maka terdapat $\delta_1$,
$\delta_2>0$ sedemikian sehingga

\Centerline{$\sum_{i=1}^n|f(b_i)-f(a_i)|\le\Bover12\epsilon$}

\noindent kapan pun
$a\le a_1\le b_1\le a_2\le b_2\le\ldots\le a_n\le b_n\le b$ dan
$\sum_{i=1}^nb_i-a_i\le\delta_1$,

\Centerline{$\sum_{i=1}^n|g(b_i)-g(a_i)|\le\Bover12\epsilon$}

\noindent kapan pun $a\le a_1\le b_1\le a_2\le b_2\le\ldots\le a_n\le
b_n\le b$ dan
$\sum_{i=1}^nb_i-a_i\le\delta_2$. Tetapkan
$\delta=\min(\delta_1,\delta_2)>0$, dan andaikan bahwa
$a\le a_1\le b_1\le a_2\le b_2\le\ldots\le a_n\le b_n\le b$ dan
$\sum_{i=1}^nb_i-a_i\le\delta$. Maka


\Centerline{$\sum_{i=1}^n|(f+g)(b_i)-(f+g)(a_i)|\le
\sum_{i=1}^n|f(b_i)-f(a_i)|+\sum_{i=1}^n|g(b_i)-g(a_i)|\le\epsilon$.}

\noindent Karena $\epsilon$ sebarang, $f+g$ kontinu mutlak.

\medskip

\quad{\bf (ii)} Misalkan $\epsilon>0$. Maka terdapat $\delta>0$ sedemikian
sehingga

\Centerline{$\sum_{i=1}^n|f(b_i)-f(a_i)|\le\Bover{\epsilon}{1+|c|}$}

\noindent kapan pun
$a\le a_1\le b_1\le a_2\le b_2\le\ldots\le a_n\le b_n\le b$ dan
$\sum_{i=1}^nb_i-a_i\le\delta$. Sekarang

\Centerline{$\sum_{i=1}^n|(cf)(b_i)-(cf)(a_i)|\le\epsilon$}

\noindent kapan pun
$a\le a_1\le b_1\le a_2\le b_2\le\ldots\le a_n\le b_n\le b$ dan
$\sum_{i=1}^nb_i-a_i\le\delta$. Karena $\epsilon$ sebarang, $cf$ kontinu
mutlak.

\medskip

{\bf (d)} Baik menurut (a) maupun menurut (b), $f$ dan $g$ terbatas; tetapkan
$M=\sup_{x\in[a,b]}|f(x)|$, $M'=\sup_{x\in[a,b]}|g(x)|$. Misalkan
$\epsilon>0$. Maka terdapat $\delta_1$, $\delta_2>0$ sedemikian sehingga

\inset{$\sum_{i=1}^n|f(b_i)-f(a_i)|\le\epsilon$
kapan pun $a\le a_1\le b_1\le a_2\le b_2\le\ldots\le a_n\le b_n\le b$ dan
$\sum_{i=1}^nb_i-a_i\le\delta_1$,}

\inset{$\sum_{i=1}^n|g(b_i)-g(a_i)|\le\epsilon$
kapan pun $a\le a_1\le b_1\le a_2\le b_2\le\ldots\le a_n\le b_n\le b$ dan
$\sum_{i=1}^nb_i-a_i\le\delta_2$.}

\noindent Tetapkan $\delta=\min(\delta_1,\delta_2)>0$ dan andaikan bahwa
$a\le a_1\le b_1\le\ldots\le b_n\le b$ serta
$\sum_{i=1}^nb_i-a_i\le\delta$. Maka

$$\eqalign{\sum_{i=1}^n|f(b_i)g(b_i)-f(a_i)g(a_i)|
&=\sum_{i=1}^n|(f(b_i)-f(a_i))g(b_i)+f(a_i)(g(b_i)-g(a_i))|\cr
&\le\sum_{i=1}^n|f(b_i)-f(a_i)||g(b_i)|+|f(a_i)||g(b_i)-g(a_i)|\cr
&\le\sum_{i=1}^n|f(b_i)-f(a_i)|M'+M|g(b_i)-g(a_i)|\cr
&\le\epsilon M'+M\epsilon
=\epsilon(M+M').\cr}$$

\noindent Karena $\epsilon$ sebarang, $f\times g$ kontinu mutlak.

\medskip

{\bf (e)} Misalkan $\epsilon>0$. Maka terdapat $\delta>0$ sedemikian sehingga
$\sum_{i=1}^n|f(d_i)-f(c_i)|\le\epsilon$ kapan pun
$c\le c_1\le d_1\le\ldots\le c_n\le d_n\le d$ dan
$\sum_{i=1}^nd_i-c_i\le\delta$; dan terdapat $\eta>0$ sedemikian sehingga
$\sum_{i=1}^n|g(b_i)-g(a_i)|\le\delta$ kapan pun
$a\le a_1\le b_1\le\ldots\le a_n\le b_n\le b$ dan
$\sum_{i=1}^nb_i-a_i\le\eta$. Sekarang andaikan bahwa
$a\le a_1\le b_1\le\ldots\le a_n\le b_n\le b$ dan
$\sum_{i=1}^nb_i-a_i\le\eta$. Karena $g$ tak-menurun, kita mempunyai
$c\le g(a_1)\le\ldots\le g(b_n)\le d$ dan
$\sum_{i=1}^ng(b_i)-g(a_i)\le\delta$, sehingga
$\sum_{i=1}^n|f(g(b_i))-f(g(a_i))|\le\epsilon$; karena $\epsilon$ sebarang,
$fg$ kontinu mutlak.
}%end of proof of 225C
\leader{225D}{Lemma} Misalkan $[a,b]$ suatu interval tertutup tak kosong
dalam $\Bbb R$ dan $f:[a,b]\to\Bbb R$ suatu fungsi kontinu mutlak yang
turunannya nol hampir di mana-mana dalam $[a,b]$.
Maka $f$ konstan pada $[a,b]$.

\proof{ Misalkan $x\in[a,b]$, $\epsilon>0$. Misalkan $\delta>0$
sedemikian sehingga $\sum_{i=1}^n|f(b_i)-f(a_i)|\le\epsilon$
kapan pun $a\le a_1\le b_1\le a_2\le b_2\le\ldots\le a_n\le b_n\le b$ dan
$\sum_{i=1}^nb_i-a_i\le\delta$. Tetapkan
$A=\{t:a<t<x,\,f'(t)$ terdefinisi dan nilainya $=0\}$; maka
$\mu A=x-a$, dengan $\mu$
menyatakan ukuran Lebesgue.
Misalkan $\Cal I$ himpunan interval tertutup tak kosong dan bukan singleton
$[c,d]\subseteq[a,x]$ sedemikian sehingga
$|f(d)-f(c)|\le\epsilon (d-c)$; maka setiap anggota $A$ termasuk dalam
anggota-anggota $\Cal I$ yang sependek apa pun.
Menurut teorema Vitali (221A), terdapat suatu subkeluarga terhitung dan
saling lepas $\Cal I_0\subseteq\Cal I$ sedemikian sehingga
$\mu(A\setminus\bigcup\Cal I_0)=0$, yakni,

\Centerline{$x-a=\mu(\bigcup\Cal I_0)=\sum_{I\in\Cal I_0}\mu I$.}

\noindent Sekarang terdapat suatu $\Cal I_1\subseteq\Cal I_0$ berhingga
sedemikian sehingga

\Centerline{$\mu(\bigcup\Cal I_1)
=\sum_{I\in\Cal I_1}\mu I\ge x-a-\delta$.}

\noindent Jika $\Cal I_1=\emptyset$, maka $x\le a+\delta$ dan
$|f(x)-f(a)|\le\epsilon$. Jika tidak, nyatakan $\Cal I_1$ sebagai
$\{[c_0,d_0],\ldots,[c_n,d_n]\}$, dengan
$a\le c_0<d_0<c_1<d_1<\ldots<c_n<d_n\le x$. Maka

\Centerline{$(c_0-a)+\sum_{i=1}^n(c_i-d_{i-1})
+(x-d_n)=\mu([a,x]\setminus\bigcup\Cal I_1)\le\delta$,}

\noindent sehingga

\Centerline{$|f(c_0)-f(a)|+\sum_{i=1}^n|f(c_i)-f(d_{i-1})|+|f(x)-f(d_n)|
\le\epsilon$.}

\noindent Di sisi lain, $|f(d_i)-f(c_i)|\le\epsilon(d_i-c_i)$ untuk
setiap $i$, sehingga

\Centerline{$\sum_{i=0}^n|f(d_i)-f(c_i)|\le\epsilon\sum_{i=0}^nd_i-c_i
\le\epsilon(x-a)$.}

\noindent Dengan menggabungkan kedua hal ini,

$$\eqalign{|f(x)-f(a)|
&\le|f(c_0)-f(a)|+|f(d_0)-f(c_0)|+|f(c_1)-f(d_0)|+\ldots\cr
&\hskip10em+|f(d_n)-f(c_n)|+|f(x)-f(d_n)|\cr
&=|f(c_0)-f(a)|+\sum_{i=1}^n|f(c_i)-f(d_{i-1})|\cr
&\hskip10em+|f(x)-f(d_n)|+\sum_{i=0}^n|f(d_i)-f(c_i)|\cr
&\le\epsilon+\epsilon(x-a)
=\epsilon(1+x-a).\cr}$$

\noindent Karena $\epsilon$ sebarang, $f(x)=f(a)$. Karena $x$ sebarang,
$f$ konstan.
}%end of proof of 225D

\leader{225E}{Teorema} Misalkan $[a,b]$ suatu interval tertutup tak kosong
dalam $\Bbb R$ dan $F:[a,b]\to\Bbb R$ suatu fungsi. Maka pernyataan-pernyataan
berikut ekuivalen:

(i) terdapat suatu fungsi bernilai real terintegralkan $f$ sedemikian sehingga
$F(x)=F(a)+\int_a^xf$ untuk setiap $x\in[a,b]$;

(ii) $\int_a^xF'$ ada dan sama dengan $F(x)-F(a)$ untuk setiap
$x\in[a,b]$;

(iii) $F$ kontinu mutlak.

\cmmnt{\medskip

\noindent{\bf Catatan} Di sini dan dalam bagian selebihnya (kecuali
dalam 225Oa), integral diambil terhadap ukuran Lebesgue pada $\Bbb R$.
}

\proof{{\bf (i)$\Rightarrow$(iii)} Andaikan (i).
Misalkan $\epsilon>0$. Menurut 225A, terdapat $\delta>0$ sedemikian sehingga
$\int_H|f|\le\epsilon$ kapan pun $H\subseteq[a,b]$ dan
$\mu H\le\delta$, dengan $\mu$ seperti biasa menyatakan ukuran Lebesgue.
Sekarang andaikan bahwa
$a\le a_1\le b_1\le a_2\le b_2\le\ldots\le a_n\le b_n\le b$ dan
$\sum_{i=1}^nb_i-a_i\le\delta$. Tinjau
$H=\bigcup_{1\le i\le n}\coint{a_i,b_i}$. Maka $\mu H\le\delta$ dan

\Centerline{$\sum_{i=1}^n|F(b_i)-F(a_i)|
=\sum_{i=1}^n|\int_{\coint{a_i,b_i}}f|
\le\sum_{i=1}^n\int_{\coint{a_i,b_i}}|f|
=\int_H|f|\le\epsilon$.}

\noindent Karena $\epsilon$ sebarang, $F$ kontinu mutlak.

\medskip

{\bf (iii)$\Rightarrow$(ii)} Jika $F$ kontinu mutlak, maka fungsi tersebut bervariasi
terbatas (225Cb), sehingga $\int_a^bF'$ ada (224I). Tetapkan
$G(x)=\int_a^xF'$ untuk $x\in [a,b]$; maka $G'\eae F'$ (222E) dan $G$
kontinu mutlak (menurut (i)$\Rightarrow$(iii) yang baru saja dibuktikan).
Dengan demikian $G-F$ kontinu mutlak (225Cc) dan terdiferensialkan hampir
di mana-mana, dengan turunan nol. Menurut 225D, $G-F$ harus konstan. Tetapi
karena $G(a)=0$, $G=F-F(a)$; tepat seperti yang disyaratkan oleh (ii).

\medskip

{\bf (ii)$\Rightarrow$(i)} langsung.
}%end of proof of 225E

\leader{225F}{Integrasi \dvrocolon{parsial}}\cmmnt{ Sebagai suatu
penerapan hasil ini, saya memberikan pembenaran bagi suatu rumus yang sudah
dikenal.

\medskip

\noindent}{\bf Teorema} Misalkan $f$ suatu fungsi bernilai real yang
terintegralkan pada suatu interval $[a,b]\subseteq\Bbb R$, dan
$g:[a,b]\to\Bbb R$ suatu fungsi kontinu mutlak. Andaikan $F$ suatu integral
tak tentu dari $f$, sehingga $F(x)-F(a)=\int_a^xf$ untuk $x\in[a,b]$. Maka

\Centerline{$\int_a^bf\times g
=F(b)g(b) - F(a)g(a) -\int_a^bF\times g'$.}

\proof{ Tetapkan $h=F\times g$. Karena $F$ kontinu mutlak (225E), demikian
pula $h$ (225Cd). Akibatnya $h(b)-h(a)=\int_a^bh'$, menurut
(iii)$\Rightarrow$(ii) dari 225E. Tetapi
$h'=F'\times g+F\times g'$ di setiap tempat di mana $F'$ dan $g'$
terdefinisi, yakni hampir di mana-mana, dan $F'\eae f$, sekali lagi menurut
222E; jadi $h'\eae f\times g+F\times g'$. Akhirnya, $g$ dan $F$ kontinu,
sehingga terukur dan terbatas, sedangkan $f$ dan $g'$ terintegralkan (dengan
kembali menggunakan 225E); jadi $f\times g$ dan $F\times g'$ terintegralkan,
dan

\Centerline{$F(b)g(b)-F(a)g(a)
=h(b)-h(a)
=\int_a^bh'
=\int_a^bf\times g+\int_a^bF\times g'$,}

\noindent seperti yang diperlukan.
}%end of proof of 225F
\leader{225G}{}\cmmnt{ Sekarang saya beralih ke sekelompok hasil pada
tingkat yang agak lebih dalam daripada sebagian besar pembahasan bab ini,
yang lebih dekat dengan gagasan-gagasan Bab 26.

\medskip

\noindent}{\bf Proposisi} Misalkan $[a,b]$ suatu interval tertutup tak kosong
dalam $\Bbb R$ dan $f:[a,b]\to\Bbb R$ suatu fungsi kontinu mutlak.

(a) $f[A]$ terabaikan untuk setiap himpunan terabaikan
$A\subseteq\Bbb R$.

(b)\dvAnew{2013}
$f[E]$ terukur untuk setiap himpunan terukur $E\subseteq\Bbb R$.

\proof{{\bf (a)} Misalkan $\epsilon>0$. Maka terdapat $\delta>0$ sedemikian
sehingga $\sum_{i=1}^n|f(b_i)-f(a_i)|\le\epsilon$ kapan pun
$a\le a_1\le b_1\le\ldots\le a_n\le b_n\le b$ dan
$\sum_{i=1}^nb_i-a_i\le\delta$. Sekarang terdapat suatu barisan
$\sequence{k}{I_k}$ dari interval tertutup yang menutupi $A$, dengan
$\sum_{k=0}^{\infty}\mu I_k\le\delta$. Untuk setiap $m\in\Bbb N$, misalkan
$F_m$ adalah $[a,b]\cap\bigcup_{k\le m}I_k$. Maka
$\mu f[F_m]\le\epsilon$. \Prf\ $F_m$ harus dapat dinyatakan sebagai
$\bigcup_{i\le n}[c_i,d_i]$, dengan $n\le m$ dan
$a\le c_0\le d_0\le\ldots\le c_n\le d_n\le b$.
Untuk setiap $i\le n$, pilih $x_i$, $y_i$ sedemikian sehingga
$c_i\le x_i$, $y_i\le d_i$ dan

\Centerline{$f(x_i)=\min_{x\in[c_i,d_i]}f(x)$,\quad
$f(y_i)=\max_{x\in[c_i,d_i]}f(x)$;}

\noindent titik-titik seperti itu ada karena $f$ kontinu (225Ca), sehingga
terbatas dan mencapai batas-batasnya pada $[c_i,d_i]$. Tetapkan
$a_i=\min(x_i,y_i)$, $b_i=\max(x_i,y_i)$, sehingga
$c_i\le a_i\le b_i\le d_i$. Maka

\Centerline{$\sum_{i=0}^nb_i-a_i\le\sum_{i=0}^nd_i-c_i=\mu F_m
\le\mu(\bigcup_{k\in\Bbb N}I_k)\le\delta$,}

\noindent sehingga

$$\eqalign{\mu f[F_m]
&=\mu(\bigcup_{i\le n}f[\,[c_i,d_i]\,])
\le\sum_{i=0}^n\mu(f[\,[c_i,d_i]\,])\cr
&=\sum_{i=0}^n\mu[f(x_i),f(y_i)]
=\sum_{i=0}^n|f(b_i)-f(a_i)|
\le\epsilon.  \text{ \Qed}\cr}$$

\noindent Tetapi $\sequence{m}{f[F_m]}$ merupakan barisan tak-menurun yang
menutupi $f[A]$, sehingga

\Centerline{$\mu^*f[A]\le\mu(\bigcup_{m\in\Bbb N}f[F_m])
=\sup_{m\in\Bbb N}\mu f[F_m]\le\epsilon$.}

\noindent Karena $\epsilon$ sebarang, $f[A]$ terabaikan, seperti yang
dinyatakan.

\medskip

{\bf (b)} Menurut 134Fb, terdapat suatu barisan $\sequencen{F_n}$ dari
subhimpunan tertutup $E\cap[a,b]$ sedemikian sehingga
$\lim_{n\to\infty}\mu F_n=\mu(E\cap[a,b])$.
Untuk setiap $n$, $F_n$ tertutup dan terbatas, jadi kompak (2A2F); karena
$f$ kontinu, $f[F_n]$ kompak (2A2Eb), jadi tertutup (2A2F, dalam arah
sebaliknya) dan terukur (114G). Selanjutnya, dengan menetapkan
$A=E\cap[a,b]\setminus\bigcup_{n\in\Bbb N}F_n$, himpunan $A$ terabaikan,
sehingga $f[A]$ terabaikan menurut (a) di sini, dan karenanya terukur.
Akibatnya

\Centerline{$f[E]=f[E\cap[a,b]]=f[\bigcup_{n\in\Bbb N}F_n\cup A]
=\bigcup_{n\in\Bbb N}f[F_n]\cup f[A]$}

\noindent terukur, seperti yang dinyatakan.
}%end of proof of 225G

\leader{225H}{Fungsi semikontinu}\cmmnt{ Sebagai persiapan bagi hasil utama
terakhir dalam bagian ini, saya memberikan suatu hasil umum mengenai fungsi
bernilai real terukur pada subhimpunan $\Bbb R$. Untuk sekali ini, akan
memudahkan jika kita meninjau fungsi-fungsi yang mengambil nilai dalam
$[-\infty,\infty]$.} Jika $D\subseteq\BbbR^r$, suatu fungsi
$g:D\to[-\infty,\infty]$ disebut {\bf semikontinu bawah} jika
$\{x:g(x)>u\}$ merupakan subhimpunan terbuka dari $D$\cmmnt{ (untuk topologi
subruang, lihat 2A3C)} untuk setiap $u\in[-\infty,\infty]$. Setiap fungsi
semikontinu bawah terukur Borel, dan karenanya terukur
Lebesgue\cmmnt{ (121B-121D)}. %121B 121C 121D
\cmmnt{Sekarang kita mempunyai hasil berikut.}

\leader{225I}{Proposisi} Andaikan $r\ge 1$ dan $f$ suatu fungsi bernilai real
yang didefinisikan pada suatu subhimpunan $D$ dari $\BbbR^r$ dan
terintegralkan pada $D$. Maka untuk setiap $\epsilon>0$ terdapat suatu fungsi
semikontinu bawah $g:\BbbR^r\to[-\infty,\infty]$ sedemikian sehingga
$g(x)\ge f(x)$ untuk setiap $x\in D$ dan $\int_Dg$ terdefinisi serta tidak
lebih besar daripada $\epsilon+\int_Df$.

\cmmnt{\medskip

\noindent{\bf Catatan} Ini merupakan hasil yang sangat penting secara umum,
jadi saya menyatakannya dalam bentuk yang cukup umum; tetapi untuk bab ini
kita hanya memerlukan kasus $r=1$, $D=[a,b]$, dengan $a\le b$.
}%end of comment

\proof{{\bf (a)} Kita dapat mengenumerasi $\Bbb Q$ sebagai
$\sequencen{q_n}$. Menurut 225A, terdapat $\delta>0$ sedemikian sehingga
$\int_F|f|\le\bover12{\epsilon}$ kapan pun $\mu_DF\le\delta$, dengan
$\mu_D$ ukuran subruang pada $D$, sehingga $\mu_DF=\mu^*F$, yaitu ukuran
luar Lebesgue dari $F$, untuk setiap $F\in\Sigma_D$, ranah $\mu_D$
(214A-214B). Untuk setiap $n\in\Bbb N$, tetapkan

\Centerline{$\delta_n=2^{-n-1}\min(\Bover{\epsilon}{1+2|q_n|},\delta)$,}

\noindent sehingga
$\sum_{n=0}^{\infty}\delta_n|q_n|\le\bover12{\epsilon}$ dan
$\sum_{n=0}^{\infty}\delta_n\le\delta$. Untuk setiap $n\in\Bbb N$,
misalkan $E_n\subseteq\BbbR^r$ suatu himpunan terukur Lebesgue sedemikian
sehingga $\{x:f(x)\ge q_n\}=D\cap E_n$, dan pilih suatu himpunan terbuka
$G_n\supseteq E_n\cap B(\tbf{0},n)$ sedemikian sehingga
$\mu G_n\le\mu(E_n\cap B(\tbf{0},n))+\delta_n$ (134Fa), dengan
$B(\tbf{0},n)$ menyatakan bola $\{x:\|x\|\le n\}$. Untuk
$x\in\BbbR^r$, tetapkan

\Centerline{$g(x)=\sup\{q_n:x\in G_n\}$,}

\noindent dengan mengizinkan $-\infty$ sebagai $\sup\emptyset$ dan
$\infty$ sebagai supremum suatu himpunan yang tidak berbatas atas dalam
$\Bbb R$.

\medskip

{\bf (b)} Sekarang periksa sifat-sifat $g$.

\medskip

\quad{\bf (i)} $g$ semikontinu bawah. \Prf\ Jika
$u\in[-\infty,\infty]$, maka

\Centerline{$\{x:g(x)>u\}=\bigcup\{G_n:q_n>u\}$}

\noindent merupakan gabungan himpunan-himpunan terbuka, jadi terbuka.\ \Qed

\medskip

\quad{\bf (ii)} $g(x)\ge f(x)$ untuk setiap $x\in D$. \Prf\ Jika
$x\in D$ dan $\eta>0$, terdapat $n\in\Bbb N$ sedemikian sehingga
$\|x\|\le n$ dan $f(x)-\eta\le q_n\le f(x)$; sekarang
$x\in E_n\subseteq G_n$, sehingga $g(x)\ge q_n\ge f(x)-\eta$. Karena
$\eta$ sebarang, $g(x)\ge f(x)$.\ \Qed

\medskip

\quad{\bf (iii)} Tinjau fungsi-fungsi
$h_1$, $h_2:D\to\ocint{-\infty,\infty}$ yang didefinisikan dengan menetapkan

$$\eqalign{h_1(x)
&=|f(x)|\text{ jika }x\in D\cap\bigcup_{n\in\Bbb N}(G_n\setminus E_n),\cr
&=0\text{ untuk }x\in D\text{ lainnya},\cr
h_2(x)&=\sum_{n=0}^{\infty}|q_n|\chi(G_n\setminus E_n)(x)
\text{ untuk setiap }x\in D.\cr}$$

\noindent Dengan menetapkan $F=\bigcup_{n\in\Bbb N}G_n\setminus E_n$,

\Centerline{$\mu F\le\sum_{n=0}^{\infty}\mu(G_n\setminus E_n)
\le\delta$,}

\noindent sehingga

\Centerline{$\int_Dh_1
=\int_{D\cap F}|f|\le\Bover12\epsilon$}

\noindent menurut pemilihan $\delta$. Adapun untuk $h_2$, kita mempunyai
(menurut teorema B. Levi)

\Centerline{$\int_Dh_2
=\sum_{n=0}^{\infty}|q_n|\mu_D(D\cap G_n\setminus E_n)
\le\sum_{n=0}^{\infty}|q_n|\mu(G_n\setminus E_n)
\le\Bover12\epsilon$}

\noindent --- karena nilai ini berhingga, $h_2(x)<\infty$ untuk hampir setiap
$x\in D$. Jadi $\int_Dh_1+h_2\le\epsilon$.

\medskip

\quad{\bf (iv)} Pokoknya ialah bahwa $g\le f+h_1+h_2$ di setiap titik dalam
$D$. \Prf\ Ambil sembarang $x\in D$. Jika $n\in\Bbb N$ dan $x\in G_n$,
maka entah $x\in E_n$, dalam hal ini

\Centerline{$f(x)+h_1(x)+h_2(x)\ge f(x)\ge q_n$,}

\noindent atau $x\in G_n\setminus E_n$, dalam hal ini

\Centerline{$f(x)+h_1(x)+h_2(x)\ge f(x)+|f(x)|+|q_n|\ge q_n$.}

\noindent Dengan demikian

\Centerline{$f(x)+h_1(x)+h_2(x)\ge\sup\{q_n:x\in G_n\}\ge g(x)$.   \Qed}

\noindent Jadi $g\le f+h_1+h_2$ di setiap titik dalam $D$.

\medskip

\quad{\bf (v)} Dengan menggabungkan (iii) dan (iv),

\Centerline{$\int_Dg\le\int_Df+h_1+h_2\le\epsilon+\int_Df$,}

\noindent seperti yang diperlukan.
}%end of proof of 225I

\leader{225J}{}\cmmnt{ Kita memerlukan beberapa hasil mengenai himpunan dan
fungsi terukur Borel yang juga menarik secara tersendiri.

\medskip

\noindent}{\bf Teorema} Misalkan $D$ suatu subhimpunan $\Bbb R$ dan
$f:D\to\Bbb R$ sembarang fungsi. Maka

\Centerline{$E=\{x:x\in D,\,f$ kontinu di $x\}$}

\noindent terukur Borel relatif dalam $D$, dan

\Centerline{$F=\{x:x\in D,\,f$ terdiferensialkan di $x\}$}

\noindent terukur Borel dalam $\Bbb R$; selain itu, $f':F\to\Bbb R$
terukur Borel.

\wheader{225J}{0}{0}{0}{36pt}

\proof{{\bf (a)} Untuk $k\in\Bbb N$, tetapkan

\Centerline{$\Cal G_k=\{\ooint{a,b}:a,\,b\in\Bbb R,\,
|f(x)-f(y)|\le 2^{-k}$ untuk semua $x$, $y\in D\cap\ooint{a,b}\}$.}

\noindent Maka $G_k=\bigcup\Cal G_k$ suatu himpunan terbuka, sehingga
$E_0=\bigcap_{k\in\Bbb N}G_k$ suatu himpunan Borel. Tetapi
$E=D\cap E_0$, jadi $E$ suatu subhimpunan Borel relatif dari $D$.

\medskip

{\bf (b)(i)} Barangkali perlu segera saya katakan bahwa ketika menafsirkan
rumus $f'(x)=\lim_{h\to 0}(f(x+h)-f(x))/h$, saya bersikeras menggunakan
definisi ketat

\Centerline{$a=\lim_{h\to 0}\bover{f(x+h)-f(x)}{h}$}

\noindent jika

\qquad\quad untuk setiap $\epsilon>0$ terdapat $\delta>0$ sedemikian sehingga
$\Bover{f(x+h)-f(x)}{h}$ terdefinisi dan

\hfill$|\Bover{f(x+h)-f(x)}{h}-a|\le\epsilon$ kapan pun
$0<|h|\le\delta$.\qquad\qquad\quad

\noindent Jadi $f'(x)$ dapat terdefinisi hanya jika terdapat
$\delta>0$ sedemikian sehingga seluruh interval
$[x-\delta,x+\delta]$ terletak dalam ranah $D$ dari $f$.

\medskip

\quad{\bf (ii)} Untuk $p$, $q$, $q'\in\Bbb Q$ dan $k\in\Bbb N$, tetapkan

$$\eqalign{H(k,p,q,q')
&=\{x:x\in E\cap\ooint{q,q'},\,|f(y)-f(x)-p(y-x)|\le 2^{-k}|y-x|
\text{ untuk setiap }y\in\ooint{q,q'}\}\cr&
\mskip300mu\text{ jika }\ooint{q,q'}\subseteq D\cr
&=\emptyset\text{ jika tidak}.\cr}$$
%display works in smallprint version

\noindent Maka
$H(k,p,q,q')=E\cap\ooint{q,q'}\cap\overline{H(k,p,q,q')}$.
\Prf\ Jika
$x\in E\cap\ooint{q,q'}\cap\overline{H(k,p,q,q')}$, terdapat suatu barisan
$\sequencen{x_n}$ dalam $H(k,p,q,q')$ yang konvergen ke $x$.
Karena $f$ kontinu di $x$,

$$\eqalign{|f(y)-f(x)-p(y-x)|
&=\lim_{n\to\infty}|f(y)-f(x_n)-p(y-x_n)|\cr&
\le 2^{-k}\lim_{n\to\infty}|y-x_n|
=2^{-k}|y-x|\cr}$$

\noindent untuk setiap $y\in\ooint{q,q'}$, sehingga
$x\in H(k,p,q,q')$.\ \QeD\ Dengan demikian $H(k,p,q,q')$ suatu himpunan
Borel. \Prf\ Menurut (a), terdapat suatu himpunan Borel $E_0$ sedemikian
sehingga $E=E_0\cap D$, sehingga jika $\ooint{q,q'}\subseteq D$, maka

\Centerline{$H(k,p,q,q')=E\cap\ooint{q,q'}\cap\overline{H(k,p,q,q')}
=E_0\cap\ooint{q,q'}\cap\overline{H(k,p,q,q')}$}

\noindent merupakan himpunan Borel. Jika tidak, tentu saja
$H(k,p,q,q')$ Borel karena kosong.\ \Qed

\medskip

\quad{\bf (iii)} Sekarang

\Centerline{$F
=\bigcap_{k\in\Bbb N}\bigcup_{p,q,q'\in\Bbb Q}H(k,p,q,q')$.}

\noindent \Prf\ ($\alpha$) Andaikan $x\in F$, yakni $f'(x)$ terdefinisi;
katakanlah $f'(x)=a$. Ambil sembarang $k\in\Bbb N$. Maka terdapat
$p\in\Bbb Q$, $\delta>0$ sedemikian sehingga
$|p-a|\le 2^{-k-1}$ dan $[x-\delta,x+\delta]\subseteq D$ serta
$|\bover{f(x+h)-f(x)}{h}-a|\le 2^{-k-1}$ kapan pun
$0<|h|\le\delta$; sekarang ambil
$q\in\Bbb Q\cap\coint{x-\delta,x}$,
$q'\in\Bbb Q\cap\ocint{x,x+\delta}$ dan lihat bahwa
$x\in H(k,p,q,q')$. Karena $x$ sebarang,
$F\subseteq\bigcap_{k\in\Bbb N}\bigcup_{p,q,q'\in\Bbb Q}H(k,p,q,q')$.
($\beta$) Jika
$x\in\bigcap_{k\in\Bbb N}\bigcup_{p,q,q'\in\Bbb Q}H(k,p,q,q')$, maka untuk
setiap $k\in\Bbb N$ pilih $p_k$, $q_k$, $q'_k\in\Bbb Q$ sedemikian sehingga
$x\in H(k,p_k,q_k,q'_k)$. Jika $h\ne 0$ dan
$x+h\in\ooint{q_k,q'_k}$, maka
$|\bover{f(x+h)-f(x)}{h}-p_k|\le 2^{-k}$.
Tetapi ini berarti, pertama, bahwa $|p_k-p_l|\le 2^{-k}+2^{-l}$ untuk setiap
$k$, $l$ (karena tentu terdapat suatu $h\ne 0$ sedemikian sehingga
$x+h\in\ooint{q_k,q'_k}\cap\ooint{q_l,q'_l}$), sehingga
$\sequence{k}{p_k}$ merupakan barisan Cauchy, katakanlah dengan limit $a$;
dan, kedua, bahwa
$|\bover{f(x+h)-f(x)}{h}-a|\le 2^{-k}+|a-p_k|$ kapan pun
$h\ne 0$ dan $x+h\in\ooint{q_k,q'_k}$, sehingga $f'(x)$ terdefinisi dan
sama dengan $a$.\ \Qed

\medskip

\quad{\bf (iv)} Karena $\Bbb Q$ terhitung, semua gabungan
$\bigcup_{p,q,q'\in\Bbb Q}H(k,p,q,q')$ merupakan himpunan Borel, sehingga
$F$ juga demikian.

\medskip

\quad{\bf (v)} Sekarang enumerasikan $\BbbQ^3$ sebagai
$\sequence{i}{(p_i,q_i,q'_i)}$, dan tetapkan
$H'_{ki}=H(k,p_i,q_i,q'_i)\setminus\bigcup_{j<i}H(k,p_j,q_j,q'_j)$ untuk
setiap $k$, $i\in\Bbb N$. Setiap $H'_{ki}$ terukur Borel,
$\sequence{i}{H'_{ki}}$ saling lepas, dan

\Centerline{$\bigcup_{i\in\Bbb N}H'_{ki}
=\bigcup_{i\in\Bbb N}H(k,p_i,q_i,q'_i)\supseteq F$}

\noindent untuk setiap $k$. Perhatikan bahwa
$|f'(x)-p|\le 2^{-k}$ kapan pun $x\in F\cap H(k,p,q,q')$, jadi jika kita
menetapkan $f_k(x)=p_i$ untuk setiap $x\in H'_{ki}$, kita memperoleh suatu
fungsi terukur Borel $f_k$ sedemikian sehingga
$|f'(x)-f_k(x)|\le 2^{-k}$ untuk setiap $x\in F$.
Dengan demikian $f'=\lim_{k\to\infty}f_k\restr F$ terukur Borel.
}%end of proof of 225J

\leader{225K}{Proposisi} Misalkan $[a,b]$ suatu interval tertutup tak kosong dalam
$\Bbb R$, dan $f:[a,b]\to\Bbb R$ suatu fungsi. Tetapkan
$F=\{x:x\in\ooint{a,b},\,f'(x)$ terdefinisi$\}$. Maka $f$ kontinu
mutlak jika dan hanya jika (i) $f$ kontinu (ii) $f'$ terintegralkan pada $F$
(iii) $f[\,[a,b]\setminus F]$ terabaikan.

\proof{{\bf (a)} Mula-mula andaikan $f$ kontinu mutlak. Maka
$f$ tentu kontinu (225Ca) dan $f'$ terintegralkan pada $[a,b]$,
karena itu juga pada $F$ (225E); selain itu, $[a,b]\setminus F$ terabaikan, sehingga
$f[\,[a,b]\setminus F]$ terabaikan, menurut 225G.

\medskip

{\bf (b)} Sekarang andaikan $f$ memenuhi syarat-syarat tersebut. Tetapkan
$f^*(x)=|f'(x)|$ untuk $x\in F$, dan $0$ untuk
$x\in[a,b]\setminus F$. Maka $f(b)\le f(a)+\int_a^bf^*$.

\medskip

\Prf\ {\bf (i)} Karena $F$ suatu himpunan Borel dan $f'$ suatu fungsi
terukur Borel (225J), $f^*$ terukur. Ambil $\epsilon>0$.
Ambil suatu subhimpunan terbuka $G$ dari
$\Bbb R$ sedemikian sehingga $f[\,[a,b]\setminus F]\subseteq G$ dan
$\mu G\le\epsilon$ (134Fa sekali lagi). Ambil suatu fungsi semikontinu bawah
$g:\Bbb R\to[0,\infty]$ sedemikian sehingga $f^*(x)\le g(x)$ untuk setiap
$x\in[a,b]$ dan $\int_a^bg\le\int_a^bf^*+\epsilon$ (225I). Perhatikan

\Centerline{$A=\{x:a\le x\le b,\,\mu([f(a),f(x)]\setminus G)
  \le 2\epsilon(x-a)+\int_a^xg\}$,}

\noindent dengan menafsirkan $[f(a),f(x)]$ sebagai $\emptyset$ jika $f(x)<f(a)$.
Maka $a\in A\subseteq[a,b]$, sehingga $c=\sup A$ terdefinisi dan termasuk dalam
$[a,b]$.

Karena $f$ kontinu, fungsi
$x\mapsto\mu([f(a),f(x)]\setminus G)$ kontinu; selain itu,
$x\mapsto 2\epsilon(x-a)+\int_a^xg$ jelas kontinu, sehingga
$c\in A$.

\medskip

\quad{\bf (ii)} \Quer\ Jika $c\in F$, sehingga $f^*(c)=|f'(c)|$, maka terdapat
suatu $\delta>0$ sedemikian sehingga

\Centerline{$a\le c-\delta\le c+\delta\le b$,}

\Centerline{$g(x)\ge g(c)-\epsilon\ge|f'(c)|-\epsilon$ kapan pun
$|x-c|\le\delta$,}

\Centerline{$|\Bover{f(x)-f(c)}{x-c}-f'(c)|\le\epsilon$ kapan pun
$0<|x-c|\le\delta$.}

\noindent Perhatikan $x=c+\delta$. Maka $c<x\le b$ dan

$$\eqalignno{\mu([f(a),f(x)]\setminus G)
&\le\mu([f(a),f(c)]\setminus G)+|f(x)-f(c)|\cr
&\le2\epsilon(c-a)+\int_a^cg+\epsilon(x-c)+|f'(c)|(x-c)\cr
&\le2\epsilon(c-a)+\int_a^cg
         +\epsilon(x-c)+\int_c^x(g+\epsilon)\cr
\noalign{\noindent (karena $g(t)\ge |f'(c)|-\epsilon$ kapan pun
$c\le t\le x$)}
&=2\epsilon(x-a)+\int_a^xg,\cr}$$

\noindent sehingga $x\in A$; tetapi $c$ seharusnya merupakan batas atas dari
$A$.\ \Bang

Jadi $c\in[a,b]\setminus F$.

\medskip

\quad{\bf (iii)} \Quer\ Sekarang andaikan, jika mungkin, bahwa
$c<b$. Kita tahu bahwa $f(c)\in G$, sehingga terdapat suatu $\eta>0$ sedemikian sehingga
$[f(c)-\eta,f(c)+\eta]\subseteq G$; sekarang terdapat suatu $\delta>0$
sedemikian sehingga $|f(x)-f(c)|\le\eta$ kapan pun $x\in[a,b]$ dan
$|x-c|\le\delta$. Tetapkan $x=\min(c+\delta,b)$; maka $c<x\le b$ dan
$[f(c),f(x)]\subseteq G$, sehingga

\Centerline{$\mu([f(a),f(x)]\setminus G)
=\mu([f(a),f(c)]\setminus G)
\le 2\epsilon(c-a)+\int_a^cg
\le 2\epsilon(x-a)+\int_a^xg$}

\noindent dan sekali lagi $x\in A$, meskipun $x>\sup A$.\ \Bang

\medskip

\quad{\bf (iv)} Kita menyimpulkan bahwa $c=b$, sehingga $b\in A$. Tetapi ini
berarti bahwa

$$\eqalign{f(b)-f(a)
&\le\mu([f(a),f(b)])
\le\mu([f(a),f(b)]\setminus G)+\mu G\cr
&\le 2\epsilon(b-a)+\int_a^bg+\epsilon
\le 2\epsilon(b-a)+\int_a^bf^*+\epsilon+\epsilon\cr
&=2\epsilon(1+b-a)+\int_a^bf^*.\cr}$$

\noindent Karena $\epsilon$ sebarang, $f(b)-f(a)\le\int_a^bf^*$, seperti
yang dinyatakan.\ \Qed

\medskip

{\bf (c)} Dengan cara yang sama, atau dengan menerapkan (b) pada $-f$,
$f(a)-f(b)\le\int_a^bf^*$, sehingga $|f(b)-f(a)|\le\int_a^bf^*$.

Tentu saja argumen tersebut berlaku pula pada setiap subinterval dari $[a,b]$, sehingga
$|f(d)-f(c)|\le\int_c^df^*$ kapan pun
$a\le c\le d\le b$. Sekarang ambil $\epsilon>0$. Menurut 225A sekali lagi, terdapat
suatu $\delta>0$ sedemikian sehingga $\int_Ef^*\le\epsilon$ kapan pun
$E\subseteq[a,b]$ dan $\mu E\le\delta$. Andaikan bahwa
$a\le a_1\le b_1\le\ldots\le a_n\le b_n\le b$ dan
$\sum_{i=1}^nb_i-a_i\le\delta$. Maka

\Centerline{$\sum_{i=1}^n|f(b_i)-f(a_i)|
\le\sum_{i=1}^n\int_{a_i}^{b_i}f^*
=\int_{\bigcupop_{i\le n}[a_i,b_i]}f^*
\le\epsilon$.}

\noindent Jadi $f$ kontinu mutlak, seperti yang dinyatakan.
}%end of proof of 225K

\leader{225L}{Korolari} Misalkan $[a,b]$ suatu interval tertutup tak kosong dalam
$\Bbb R$. Misalkan $f:[a,b]\to\Bbb R$ suatu fungsi kontinu yang
terdiferensialkan pada interval terbuka $\ooint{a,b}$. Jika turunannya
$f'$ terintegralkan pada $[a,b]$, maka $f$ kontinu mutlak, dan
$f(b)-f(a)=\int_a^bf'$.

\proof{ $f[\,[a,b]\setminus F]=\{f(a),f(b)\}$ tentu terabaikan, sehingga
$f$ kontinu mutlak, menurut 225K; akibatnya
$f(b)-f(a)=\int_a^bf'$, menurut 225E.
}%end of proof of 225L

\leader{225M}{Korolari} Misalkan $[a,b]$ suatu interval tertutup tak kosong dalam
$\Bbb R$, dan $f:[a,b]\to\Bbb R$ suatu fungsi kontinu. Maka $f$
kontinu mutlak jika dan hanya jika fungsi tersebut kontinu dan bervariasi terbatas serta
$f[A]$ terabaikan untuk setiap himpunan terabaikan $A\subseteq[a,b]$.

\proof{{\bf (a)} Andaikan bahwa $f$ kontinu mutlak. Menurut
225C(a-b), fungsi tersebut kontinu dan bervariasi terbatas, dan menurut 225G kita mempunyai
$f[A]$ terabaikan untuk setiap himpunan terabaikan $A\subseteq [a,b]$.

\medskip

{\bf (b)} Sekarang andaikan $f$ memenuhi syarat-syarat tersebut. Tetapkan
$F=\{x:x\in\ooint{a,b},\,f'(x)$ terdefinisi$\}$. Menurut 224I sekali lagi,
$[a,b]\setminus F$ terabaikan, sehingga $f[\,[a,b]\setminus F]$
terabaikan. Selain itu, juga menurut 224I, $f'$ terintegralkan pada $[a,b]$
atau $F$. Jadi syarat-syarat 225K terpenuhi dan $f$ kontinu
mutlak.
}%end of proof of 225M

\leader{225N}{Fungsi Cantor}\cmmnt{ Saya perlu menyebutkan
contoh baku suatu fungsi kontinu yang bervariasi terbatas tetapi
tidak kontinu mutlak.} Misalkan $C\subseteq[0,1]$ himpunan
Cantor\cmmnt{ (134G)}.
Ingat bahwa `fungsi Cantor' adalah suatu fungsi kontinu tak-menurun
$f:[0,1]\to[0,1]$ sedemikian sehingga $f'(x)$ terdefinisi dan sama dengan
nol untuk setiap $x\in[0,1]\setminus C$,
tetapi $f(0)=0<1=f(1)$\cmmnt{ (134H)}. \cmmnt{Tentu
saja} $f$ bervariasi terbatas dan tidak kontinu mutlak. $C$
terabaikan sedangkan $f[C]=[0,1]$ tidak. Jika $x\in C$, maka untuk setiap
$n\in\Bbb N$ terdapat suatu interval dengan panjang $3^{-n}$, yang memuat $x$, dan yang
padanya $f$ meningkat sebesar $2^{-n}$; jadi $f$ tidak mungkin terdiferensialkan di
$x$, dan himpunan $F=\dom f'$ dari 225K tepat sama dengan $[0,1]\setminus C$, sehingga
$f[\,[0,1]\setminus F]=[0,1]$.

\vleader{48pt}{225O}{Fungsi bernilai kompleks}\cmmnt{ Seperti biasa, saya
jabarkan
hasil-hasil di atas dalam bentuk yang berlaku bagi fungsi bernilai kompleks.

\medskip

} {\bf (a)} Misalkan
$(X,\Sigma,\mu)$ sembarang ruang ukur dan $f$ suatu fungsi
bernilai kompleks terintegralkan
yang didefinisikan pada suatu subhimpunan koterabaikan dari $X$.
Maka untuk setiap $\epsilon>0$ terdapat suatu himpunan terukur
$E$ berukuran hingga dan suatu bilangan real
$\delta>0$ sedemikian sehingga $\int_F|f|\le\epsilon$ kapan pun $F\in\Sigma$ dan
$\mu(F\cap E)\le\delta$. \prooflet{(Terapkan 225A pada $|f|$.)}

\spheader 225Ob Jika $[a,b]$ suatu interval tertutup
tak kosong dalam $\Bbb R$ dan $f:[a,b]\to\Bbb C$ suatu fungsi, kita katakan
bahwa $f$ {\bf kontinu mutlak} jika untuk setiap $\epsilon>0$ terdapat
suatu $\delta>0$ sedemikian sehingga $\sum_{i=1}^n|f(b_i)-f(a_i)|\le\epsilon$
kapan pun $a\le a_1\le b_1\le a_2\le b_2\le\ldots\le a_n\le b_n\le b$ dan
$\sum_{i=1}^nb_i-a_i\le\delta$. Perhatikan bahwa $f$ kontinu
mutlak jika dan hanya jika bagian real dan bagian imajinernya sama-sama
kontinu mutlak.

\spheader 225Oc Misalkan $[a,b]$ suatu interval tertutup tak kosong dalam
$\Bbb R$.

\quad(i) Jika $f:[a,b]\to\Bbb C$ kontinu mutlak, maka fungsi tersebut bervariasi
terbatas pada $[a,b]$, sehingga terdiferensialkan hampir di mana-mana dalam $[a,b]$,
dan turunannya terintegralkan pada $[a,b]$.

\quad(ii) Jika $f$, $g:[a,b]\to\Bbb C$ kontinu mutlak, maka demikian pula
$f+g$ dan $\zeta f$, untuk setiap $\zeta\in\Bbb C$, serta $f\times g$.

\quad(iii) Jika $g:[a,b]\to[c,d]$ monoton dan kontinu mutlak,
dan
$f:[c,d]\to\Bbb C$ kontinu mutlak, maka $fg:[a,b]\to\Bbb C$
kontinu mutlak.

\spheader 225Od Misalkan $[a,b]$ suatu interval tertutup tak kosong dalam
$\Bbb R$ dan $F:[a,b]\to\Bbb C$ suatu fungsi. Maka pernyataan-pernyataan berikut
ekuivalen:

\quad(i) terdapat suatu fungsi bernilai kompleks terintegralkan $f$ sedemikian sehingga
$F(x)=F(a)+\int_a^xf$ untuk setiap $x\in[a,b]$;

\quad(ii) $\int_a^xF'$ ada dan sama dengan
$F(x)-F(a)$ untuk setiap $x\in[a,b]$;

\quad(iii) $F$ kontinu mutlak.

\cmmnt{\noindent (Terapkan 225E pada bagian real dan bagian imajiner dari $F$.)}

\spheader 225Oe Misalkan $f$ suatu fungsi bernilai kompleks terintegralkan
pada suatu interval $[a,b]\subseteq\Bbb R$, dan $g:[a,b]\to\Bbb C$ suatu fungsi
kontinu mutlak. Tetapkan $F(x)=\int_a^xf$ untuk $x\in[a,b]$.
Maka

\Centerline{$\int_a^bf\times g
=F(b)g(b) - F(a)g(a) -\int_a^bF\times g'$.}

\cmmnt{\noindent (Terapkan 225F pada bagian real dan bagian imajiner dari $f$
dan $g$.)}

\spheader 225Of Misalkan $f$ suatu fungsi kontinu bernilai kompleks
pada suatu interval tertutup $[a,b]\subseteq\Bbb R$, dan andaikan $f$
terdiferensialkan di setiap titik dari interval terbuka $\ooint{a,b}$, dengan
$f'$ terintegralkan pada $[a,b]$. Maka $f$ kontinu mutlak.
\cmmnt{(Terapkan 225L pada bagian real dan bagian imajiner dari $f$.)}

\cmmnt{\spheader 225Og Untuk hasil yang bersesuaian dengan 225M, lihat 264Yp.}

\exercises{
\leader{225X}{Latihan dasar (a)}
%\spheader 225Xa
Tunjukkan langsung dari definisi dalam
225B (tanpa menggunakan
225E) bahwa setiap fungsi bernilai real yang kontinu mutlak pada suatu interval
tertutup $[a,b]$ dapat dinyatakan sebagai selisih fungsi-fungsi kontinu mutlak
yang tak-menurun.
%225B

\spheader 225Xb Tunjukkan langsung dari definisi dalam 225B dan
Teorema Nilai Rata-Rata (tanpa menggunakan 225K) bahwa jika suatu fungsi
$f$ kontinu pada interval tertutup $[a,b]$, terdiferensialkan pada interval
terbuka $\ooint{a,b}$, dan turunannya terbatas pada
$\ooint{a,b}$, maka $f$ kontinu mutlak, sehingga
$f(x)=f(a)+\int_a^xf'$ untuk setiap $x\in[a,b]$.
%225E

\spheader 225Xc Tunjukkan bahwa jika $f:[a,b]\to\Bbb R$ kontinu
mutlak, maka $\Var f=\int_a^b|f'|$.   \Hint{gabungkan 224I dan
225E.}
%225E

\spheader 225Xd Misalkan $g:\Bbb R\to\Bbb R$ suatu fungsi
tak-menurun yang kontinu mutlak pada setiap interval terbatas;
misalkan $\mu_g$ ukuran Lebesgue-Stieltjes yang bersesuaian
(114Xa), dan $\Sigma_g$ domainnya.   Tunjukkan bahwa $\int_Eg'=\mu_gE$ untuk
setiap $E\in\Sigma_g$, jika kita mengizinkan $\infty$ sebagai nilai
integral.   \Hint{mulailah dengan interval $E$.}
%225E

\spheader 225Xe Misalkan $g:[a,b]\to\Bbb R$ suatu fungsi kontinu
mutlak yang tak-menurun, dan $f:[g(a),g(b)]\to\Bbb R$ suatu fungsi
kontinu.   Tunjukkan bahwa $\int_{g(a)}^{g(b)}f(t)dt=
\int_a^bf(g(t))g'(t)dt$.   \Hint{tetapkan $F(x)=\int_{g(a)}^xf$,
$G=Fg$ dan tinjau $\int_a^bG'(t)dt$.   Lihat juga 263J.}
%225E

\spheader 225Xf Andaikan bahwa $I\subseteq\Bbb R$ sembarang
interval tak-sepele (terbatas atau tak terbatas, terbuka, tertutup, atau
setengah terbuka, tetapi bukan himpunan kosong ataupun himpunan satu anggota),
dan $f:I\to\Bbb R$ suatu fungsi.   Tunjukkan bahwa
$f$ kontinu mutlak pada setiap subinterval tertutup dan terbatas dari $I$
jika dan hanya jika terdapat fungsi $g$ sedemikian sehingga
$\int_a^bg=f(b)-f(a)$ setiap kali
$a\le b$ dalam $I$, dan dalam hal ini $g$ terintegralkan jika dan hanya jika
$f$ mempunyai variasi terbatas pada $I$.
%225E

\spheader 225Xg Tunjukkan bahwa
$\biggerint_0^1\Bover{\ln x}{x-1}dx=\sum_{n=1}^{\infty}\Bover1{n^2}$.
\Hint{gunakan 225F untuk menentukan $\int_0^1x^n\ln x\,dx$, dan ingat bahwa
$\bover1{1-x}=\sum_{n=0}^{\infty}x^n$ untuk $0\le x<1$.}
%225F

\spheader 225Xh(i) Tunjukkan bahwa $\int_0^1t^adt$ berhingga untuk
setiap $a>-1$.   (ii) Tunjukkan bahwa $\int_1^{\infty}t^ae^{-t}dt$ berhingga
untuk setiap $a\in\Bbb R$.   \Hint{tunjukkan bahwa terdapat $M$
sedemikian sehingga $t^a\le Me^{t/2}$ untuk $t\ge 1$.}   (iii) Tunjukkan bahwa
$\Gamma(a)=\int_0^{\infty}t^{a-1}e^{-t}dt$ terdefinisi untuk setiap $a>0$.
(iv) Tunjukkan bahwa $\Gamma(a+1)=a\Gamma(a)$ untuk setiap $a>0$.   (v) Tunjukkan bahwa
$\Gamma(n+1)=n!$ untuk setiap $n\in\Bbb N$.

($\Gamma$ tentu saja merupakan {\bf fungsi gama}.)
%225F

\spheader 225Xi Tunjukkan bahwa jika $b>0$, maka
$\int_0^{\infty}u^{b-1}e^{-u^2/2}du=2^{(b-2)/2}\Gamma(\bover{b}2)$.
\Hint{tinjau $f(t)=t^{(b-2)/2}e^{-t}$, $g(u)=u^2/2$ dalam 225Xe.}
%225Xe, 225Xh, 225F

\spheader 225Xj Andaikan bahwa $f$, $g$ adalah
fungsi-fungsi semikontinu bawah yang didefinisikan pada subhimpunan-subhimpunan $\BbbR^r$,
dan bernilai dalam $\ocint{-\infty,\infty}$.   (i) Tunjukkan bahwa $f+g$,
$f\wedge g$ dan $f\vee g$ semikontinu bawah, dan bahwa
$\alpha f$ semikontinu
bawah untuk setiap $\alpha\ge 0$.   (ii) Tunjukkan bahwa jika $f$ dan $g$
nonnegatif, maka $f\times g$ semikontinu bawah.
(iii) Tunjukkan bahwa jika $f$ nonnegatif dan $g$ kontinu, maka
$f\times g$ semikontinu bawah.   (iv) Tunjukkan bahwa jika $f$
tak-menurun, maka komposisi $fg$ semikontinu bawah.
%225H

\spheader 225Xk Misalkan $A$ suatu keluarga tak kosong dari fungsi-fungsi
semikontinu bawah yang didefinisikan pada subhimpunan-subhimpunan $\BbbR^r$ dan bernilai
dalam $[-\infty,\infty]$.   Tetapkan $g(x)=\sup\{f(x):f\in A,\,x\in\dom
f\}$ untuk $x\in
D=\bigcup_{f\in A}\dom f$.   Tunjukkan bahwa $g$ semikontinu
bawah.
%225H

\spheader 225Xl Misalkan $f:[a,b]\to\Bbb R$ suatu fungsi kontinu mutlak,
dengan $a\le b$.   (i) Tunjukkan bahwa $|f|:[a,b]\to\Bbb R$ kontinu
mutlak.   (ii) Tunjukkan bahwa $gf$ kontinu mutlak
setiap kali $g:f[\,[a,b]\,]\to\Bbb R$ kontinu mutlak dan
$g'$ terbatas.   (iii) Tunjukkan bahwa jika $g:[a,b]\to\Bbb R$ kontinu
mutlak dan $\inf_{x\in[a,b]}|g(x)|>0$, maka
$f/g$ kontinu mutlak.
%225K

\spheader 225Xm Andaikan bahwa $f:[a,b]\to\Bbb R$ kontinu, dan
terdiferensialkan pada semua kecuali terhitung banyak
titik dari $[a,b]$. Tunjukkan bahwa $f$ kontinu mutlak jika dan hanya jika
fungsi tersebut mempunyai variasi terbatas.
%225K

\spheader 225Xn Misalkan $f:\coint{0,\infty}\to\Bbb C$ suatu fungsi
yang kontinu mutlak pada $[0,a]$ untuk setiap
$a\in\coint{0,\infty}$ dan mempunyai transformasi Laplace
$F(s)=\int_0^{\infty}e^{-sx}f(x)dx$ yang terdefinisi pada $\{s:\Real s>S\}$.
Andaikan pula bahwa $\lim_{x\to\infty}e^{-Sx}f(x)=0$.   Tunjukkan
bahwa $f'$ mempunyai transformasi Laplace $sF(s)-f(0)$ yang terdefinisi setiap kali $\Real
s>S$.   ({\it Petunjuk\/}:  tunjukkan bahwa

\Centerline{$f(x)e^{-sx}-f(0)
=\int_0^x\Bover{d}{dt}(f(t)e^{-st})dt$}

\noindent untuk setiap $x\ge 0$.)
%225O 225F

\leader{225Y}{Latihan lanjutan (a)}
%\spheader 225Ya
Tunjukkan bahwa komposisi dua fungsi kontinu
mutlak tidak harus kontinu mutlak.   \Hint{224Xb.}
%225B

\spheader 225Yb Misalkan $f:[a,b]\to\Bbb R$ suatu fungsi kontinu, dengan
$a<b$.   Tetapkan $G=\{x:x\in\ooint{a,b},\,\exists\,y\in\ocint{x,b}$ sedemikian
sehingga $f(x)<f(y)\}$.   Tunjukkan bahwa $G$ terbuka dan dapat dinyatakan sebagai
gabungan saling lepas dari interval-interval $\ooint{c,d}$ dengan $f(c)\le f(d)$.   Gunakan
hal ini untuk membuktikan 225D tanpa menggunakan Teorema Vitali.
%225D 2A2I

\spheader 225Yc Misalkan $f:[a,b]\to\Bbb R$ suatu fungsi dengan variasi
terbatas dan $\gamma>0$.   Tunjukkan bahwa terdapat fungsi kontinu mutlak
$g:[a,b]\to\Bbb R$ sedemikian sehingga $|g'(x)|\le\gamma$ di setiap titik tempat
turunannya ada dan $\{x:x\in[a,b],\,f(x)\ne g(x)\}$ mempunyai ukuran
paling besar $\bover1{\gamma}\Var f$.   \Hint{cukup buktikan kasus
$f$ tak-menurun.   Terapkan 225Yb pada fungsi
$x\mapsto f(x)-\gamma x$ dan tunjukkan bahwa $\gamma\mu G\le\Var f$.
Tetapkan $g(x)=f(x)$ untuk $x\in\ooint{a,b}\setminus G$.}
%225Yb 225D

\spheader 225Yd\dvAnew{2013}
Misalkan $f:\Bbb R\to\Bbb R$ suatu fungsi yang kontinu
mutlak pada setiap interval terbatas.   Tunjukkan bahwa
$\Var f\le\bover12\Var f'+\int|f|$.
%225E mt22bits

\spheader 225Ye Misalkan $f$ suatu fungsi bernilai real terukur dan
nonnegatif yang didefinisikan pada suatu subhimpunan $D$ dari $\BbbR^r$, dengan $r\ge 1$.   Tunjukkan
bahwa untuk setiap $\epsilon>0$ terdapat fungsi semikontinu bawah
$g:\BbbR^r\to[-\infty,\infty]$ sedemikian sehingga $g(x)\ge f(x)$ untuk
setiap $x\in D$ dan $\int_Dg-f\le\epsilon$.
%225I

\spheader 225Yf Misalkan $f$ suatu fungsi bernilai real terukur
yang didefinisikan pada suatu subhimpunan $D$ dari $\BbbR^r$, dengan $r\ge 1$.   Tunjukkan bahwa untuk
setiap $\epsilon>0$ terdapat fungsi semikontinu bawah
$g:\BbbR^r\to[-\infty,\infty]$ sedemikian sehingga $g(x)\ge f(x)$ untuk setiap
$x\in D$ dan
$\mu^*\{x:x\in D,\,g(x)>f(x)\}\le\epsilon$.   \Hint{134Yd, 134Fb.}
%225I

\spheader 225Yg(i) Tunjukkan bahwa jika $f$ suatu fungsi real terukur Lebesgue,
maka semua turunan Dini-nya terukur Lebesgue.
(ii) Tunjukkan bahwa jika $f$ suatu fungsi real terukur Borel,
maka semua turunan Dini-nya terukur Borel.
%225J mt22bits
}%end of exercises

\endnotes{
\Notesheader{225} Masih banyak yang dapat dikatakan tentang fungsi-fungsi kontinu
mutlak; saya akan kembali ke topik ini dalam bagian berikutnya
dan dalam Bab 26.   Saya jarang akan menggunakan secara langsung hasil-hasil
mulai 225H dalam bentuknya yang paling kuat,
tetapi menurut saya penyelidikan semacam ini
diperlukan untuk memperoleh gambaran yang jelas mengenai hubungan antara
konsep-konsep seperti kontinuitas mutlak dan variasi terbatas. Tentu saja,
untuk menerapkan hasil-hasil ini, kita memang memerlukan bekal kelas-kelas
sederhana fungsi kontinu mutlak; fungsi-fungsi terdiferensialkan dengan
turunan terbatas membentuk kelas terpenting (225Xb).   Keluarga yang lebih
besar dari jenis yang sama adalah kelas fungsi `Lipschitz'
(262Bc).

Definisi `fungsi kontinu mutlak' biasanya diberikan
untuk interval tertutup dan terbatas, seperti dalam 225B.   Pokoknya ialah bahwa untuk
interval lain, generalisasi paling sederhana dari perumusan ini tidak
tampak sepenuhnya tepat.   Dalam 225Xf saya mencoba menunjukkan jenis persyaratan
yang mungkin dikenakan pada fungsi-fungsi yang didefinisikan pada jenis interval lain.

Perlu saya katakan bahwa sasaran sesungguhnya masih belum sepenuhnya dalam jangkauan kita.
Saya telah dapat memberikan perumusan yang cukup memuaskan mengenai integrasi
parsial sederhana (225F), setidaknya untuk interval terbatas -- suatu proses
limit lebih lanjut diperlukan untuk menangani interval tak terbatas.   Namun suatu
metode pendamping dari kalkulus lanjut, integrasi dengan substitusi,
masih sukar dicapai.   Hasil terbaik yang menurut saya dapat kita peroleh pada tahap ini adalah 225Xe,
yang mensyaratkan integran $f$ kontinu.   Memang benar bahwa
hasil tersebut berlaku untuk setiap $f$ yang terintegralkan, tetapi masih ada beberapa
kehalusan lebih lanjut yang harus dikuasai dalam prosesnya; gagasan-gagasan yang diperlukan diberikan dalam
hasil-hasil yang jauh lebih umum 235A dan 263D di bawah, dan diterapkan pada
kasus satu dimensi dalam 263J.

Dalam perjalanan menuju karakterisasi fungsi-fungsi kontinu mutlak dalam
225K, saya mendapati diri saya menggunakan salah satu hubungan mendasar
antara ukuran Lebesgue dan topologi $\BbbR^r$ (225I).   Teknik
di sini dapat diadaptasi untuk memberikan banyak variasi hasil tersebut;
lihat 225Ye-225Yf.   Jika Anda belum pernah menjumpai fungsi-fungsi semikontinu
sebelumnya, 225Xj-225Xk memberikan sebagian gambaran tentang sifat-sifatnya.   Dalam 225J
saya memberikan sebuah fragmen `teori himpunan deskriptif', yakni kajian tentang jenis-jenis
himpunan yang dapat muncul dari rumus-rumus analisis.   Gagasan-gagasan ini juga
akan muncul kembali di tempat lain (bandingkan 225Yg dan juga pembuktian 262M
di bawah) dan akan sangat penting dalam Jilid 4 dan 5.
}%end of notes

\discrpage
